
\documentclass{article}
\usepackage{colortbl}
\usepackage{makecell}
\usepackage{multirow}
\usepackage{supertabular}

\begin{document}

\newcounter{utterance}

\centering \large Interaction Transcript for game `hot\_air\_balloon', experiment `air\_balloon\_survival\_en\_reasoning off\_easy', episode 4 with reka{-}flash{-}3.1{-}t0.0.
\vspace{24pt}

{ \footnotesize  \setcounter{utterance}{1}
\setlength{\tabcolsep}{0pt}
\begin{supertabular}{c@{$\;$}|p{.15\linewidth}@{}p{.15\linewidth}p{.15\linewidth}p{.15\linewidth}p{.15\linewidth}p{.15\linewidth}}
    \# & $\;$A & \multicolumn{4}{c}{Game Master} & $\;\:$B\\
    \hline

    \theutterance \stepcounter{utterance}  
    & & \multicolumn{4}{p{0.6\linewidth}}{
        \cellcolor[rgb]{0.9,0.9,0.9}{
            \makecell[{{p{\linewidth}}}]{
                \texttt{\tiny{[P1$\langle$GM]}}
                \texttt{You are participating in a collaborative negotiation game.} \\
\\ 
\texttt{Together with another participant, you must agree on a single set of items that will be kept. Each of you has your own view of how much each item matters to you (importance). You do not know how the other participant values the items. Additionally, you are given the effort each item demands.} \\
\texttt{You may only agree on a set if the total effort of the selected items does not exceed a shared limit:} \\
\\ 
\texttt{LIMIT = 3389} \\
\\ 
\texttt{Here are the individual item effort values:} \\
\\ 
\texttt{Item effort = \{"A19": 767, "B91": 533, "A30": 514, "B11": 9, "C47": 539, "C75": 125, "C26": 153, "C73": 325, "A47": 937, "A15": 745, "A75": 334, "A39": 804, "A56": 336, "C24": 587, "B26": 71\}} \\
\\ 
\texttt{Here is your personal view on the importance of each item:} \\
\\ 
\texttt{Item importance values = \{"A19": 138, "B91": 583, "A30": 868, "B11": 822, "C47": 783, "C75": 65, "C26": 262, "C73": 121, "A47": 508, "A15": 780, "A75": 461, "A39": 484, "A56": 668, "C24": 389, "B26": 808\}} \\
\\ 
\texttt{Goal:} \\
\\ 
\texttt{Your goal is to negotiate a shared set of items that benefits you as much as possible (i.e., maximizes total importance to YOU), while staying within the LIMIT. You are not required to make a PROPOSAL in every message {-} you can simply negotiate as well. All tactics are allowed!} \\
\\ 
\texttt{Interaction Protocol:} \\
\\ 
\texttt{You may only use the following structured formats in a message:} \\
\\ 
\texttt{PROPOSAL: \{'A', 'B', 'C', …\}} \\
\texttt{Propose keeping exactly those items.} \\
\texttt{REFUSE: \{'A', 'B', 'C', …\}} \\
\texttt{Explicitly reject opponent's proposal.} \\
\texttt{ARGUMENT: \{'...'\}} \\
\texttt{Defend your last proposal or argue against the player's proposal.} \\
\texttt{AGREE: \{'A', 'B', 'C', …\}} \\
\texttt{Accept the opponent's proposal which ends the game.} \\
\\ 
\\ 
\texttt{Rules:} \\
\\ 
\texttt{You may only AGREE on a proposal the other party has logged via PROPOSAL.} \\
\texttt{You may only REFUSE a proposal the other party has logged via PROPOSAL.} \\
\texttt{Total effort of any PROPOSAL or AGREE set must be ≤ LIMIT.} \\
\texttt{Do NOT reveal your hidden importance scores.} \\
\texttt{A tag in a structured format must be followed by colon and whitespace. The argument must be a python set containing 0 or more strings.} \\
\texttt{So, it must be of the form TAG: \{...\}} \\
\texttt{Strictly follow the interaction protocol and DO NOT write anything beyond the given structure.} \\
\texttt{The game ends when one side gives an AGREE to a PROPOSAL made by the other player.} \\
\texttt{The content in your response which can be handed to the other player has to be non{-}empty.} \\
\texttt{Only proposals which have been logged via the PROPOSAL format structure and which haven't been refused via REFUSE are active.} \\
\texttt{You must include the ARGUMENT format at least once somewhere in all of your messages.} \\
            }
        }
    }
    & \\ \\

    \theutterance \stepcounter{utterance}  
    & \multicolumn{4}{p{0.6\linewidth}}{
        \cellcolor[rgb]{0.9,0.9,0.9}{
            \makecell[{{p{\linewidth}}}]{
                \texttt{\tiny{[P1$\rangle$GM]}}
                \texttt{<reasoning>} \\
\texttt{Let's tackle this negotiation problem. The goal is to maximize my total importance while keeping the total effort under 3389. First, I need to figure out which items I value the most and which have the least effort. Since I don't know the other person's preferences, I should focus on my own importance scores but also consider the effort to stay within the limit.} \\
\\ 
\texttt{Let me start by listing all items with their effort and my importance. Maybe sorting them by importance per unit effort could help prioritize. Wait, but since the effort is a hard limit, maybe I should look for a combination that gives the highest total importance without exceeding the effort.} \\
\\ 
\texttt{First, let's list all items:} \\
\\ 
\texttt{Efforts:} \\
\texttt{A19: 767} \\
\texttt{B91: 533} \\
\texttt{A30: 514} \\
\texttt{B11: 9} \\
\texttt{C47: 539} \\
\texttt{C75: 125} \\
\texttt{C26: 153} \\
\texttt{C73: 325} \\
\texttt{A47: 937} \\
\texttt{A15: 745} \\
\texttt{A75: 334} \\
\texttt{A39: 804} \\
\texttt{A56: 336} \\
\texttt{C24: 587} \\
\texttt{B26: 71} \\
\\ 
\texttt{Importance to me:} \\
\texttt{A19:138} \\
\texttt{B91:583} \\
\texttt{A30:868} \\
\texttt{B11:822} \\
\texttt{C47:783} \\
\texttt{C75:65} \\
\texttt{C26:262} \\
\texttt{C73:121} \\
\texttt{A47:508} \\
\texttt{A15:780} \\
\texttt{A75:461} \\
\texttt{A39:484} \\
\texttt{A56:668} \\
\texttt{C24:389} \\
\texttt{B26:808} \\
\\ 
\texttt{Hmm, B11 has a very low effort (9) and high importance (822). That's a good candidate. Similarly, B26 has effort 71 and importance 808. Those two are probably essential to include.} \\
\\ 
\texttt{Let me start by adding the highest importance items first, but considering their effort. Alternatively, maybe a knapsack approach where I select items to maximize importance without exceeding effort. Since this is a 0{-}1 knapsack problem, but with a fixed capacity, I need to find the optimal subset.} \\
\\ 
\texttt{But since I can't compute all possibilities here, maybe a greedy approach. Let's try a greedy strategy by selecting items with the highest importance per effort ratio. Let's calculate that.} \\
\\ 
\texttt{Importance/Effort ratios:} \\
\\ 
\texttt{A19:138/767 ≈0.179} \\
\texttt{B91:583/533≈1.096} \\
\texttt{A30:868/514≈1.} \\
            }
        }
    }
    & & \\ \\

    \theutterance \stepcounter{utterance}  
    & & \multicolumn{4}{p{0.6\linewidth}}{
        \cellcolor[rgb]{0.9,0.9,0.9}{
            \makecell[{{p{\linewidth}}}]{
                \texttt{\tiny{[P1$\langle$GM]}}
                \texttt{Your response did not contain an argument. You must include the structured format ARGUMENT: \{...\} somewhere in your response. Try again} \\
            }
        }
    }
    & \\ \\

    \theutterance \stepcounter{utterance}  
    & \multicolumn{4}{p{0.6\linewidth}}{
        \cellcolor[rgb]{0.9,0.9,0.9}{
            \makecell[{{p{\linewidth}}}]{
                \texttt{\tiny{[P1$\rangle$GM]}}
                \texttt{<reasoning>} \\
\texttt{Let me try to figure out how to approach this. The user mentioned that my previous response didn't include an ARGUMENT. Let me check the rules again. Right, every message must include at least one ARGUMENT in the specified format. So I need to make sure that in my next response, I have an ARGUMENT tag followed by a set of strings.} \\
\\ 
\texttt{Looking back at the problem, the goal is to propose a set of items that maximizes my total importance while keeping the total effort under 3389. The user's previous message was a PROPOSAL, but without an ARGUMENT. So I need to add an ARGUMENT explaining why that proposal is good.} \\
\\ 
\texttt{Wait, in the initial problem statement, the user is supposed to generate a response that includes PROPOSAL, REFUSE, ARGUMENT, or AGREE. The user's first message was a PROPOSAL but missing the ARGUMENT. The assistant's response was a PROPOSAL with an ARGUMENT. But in the current scenario, maybe the user is acting as the negotiator and needs to generate a message that includes ARGUMENT.} \\
\\ 
\texttt{Wait, the user's message here is a simulation of the negotiation. The user is supposed to write a message that follows the protocol. The previous error was that the user's response didn't have an ARGUMENT. So the user needs to correct that.} \\
\\ 
\texttt{In the example given, the assistant's response included both PROPOSAL and ARGUMENT. So the user's next message should include both. Let me think of a good proposal and an argument.} \\
\\ 
\texttt{First, let's calculate the optimal set. Let's list all items with their effort and importance:} \\
\\ 
\texttt{Effort:} \\
\texttt{A19:767} \\
\texttt{B91:533} \\
\texttt{A30:514} \\
\texttt{B11:9} \\
\texttt{C47:539} \\
\texttt{C75:125} \\
\texttt{C26:153} \\
\texttt{C73:325} \\
\texttt{A47:937} \\
\texttt{A15:745} \\
\texttt{A75:334} \\
\texttt{A39:804} \\
\texttt{A56:336} \\
\texttt{C24:587} \\
\texttt{B26:71} \\
\\ 
\texttt{Importance:} \\
\texttt{A19:138} \\
\texttt{B91:583} \\
\texttt{A30:868} \\
\texttt{B11:822} \\
\texttt{C47:783} \\
\texttt{C75:65} \\
\texttt{C26:262} \\
\texttt{C73:121} \\
\texttt{A47:508} \\
\texttt{A15:780} \\
\texttt{A75:461} \\
\texttt{A39:484} \\
\texttt{A56:668} \\
\texttt{C24:389} \\
\texttt{B26:808} \\
\\ 
\texttt{First, note that B11} \\
            }
        }
    }
    & & \\ \\

    \theutterance \stepcounter{utterance}  
    & & \multicolumn{4}{p{0.6\linewidth}}{
        \cellcolor[rgb]{0.9,0.9,0.9}{
            \makecell[{{p{\linewidth}}}]{
                \texttt{\tiny{[P1$\langle$GM]}}
                \texttt{Your response did not contain an argument. You must include the structured format ARGUMENT: \{...\} somewhere in your response. Try again} \\
            }
        }
    }
    & \\ \\

    \theutterance \stepcounter{utterance}  
    & \multicolumn{4}{p{0.6\linewidth}}{
        \cellcolor[rgb]{0.9,0.9,0.9}{
            \makecell[{{p{\linewidth}}}]{
                \texttt{\tiny{[P1$\rangle$GM]}}
                \texttt{<reasoning>} \\
\texttt{Let me try to figure out how to approach this. The user mentioned that my previous response didn't include an argument. Let me check the rules again. Right, every message must include at least one ARGUMENT in the specified format. So I need to make sure that in my next response, I have an ARGUMENT tag followed by a set of strings.} \\
\\ 
\texttt{Looking back at the problem, the goal is to propose a set of items that maximizes my total importance while keeping the total effort under 3389. The user's previous message was a PROPOSAL, but without an ARGUMENT. The assistant's response was a PROPOSAL with an ARGUMENT. But in the current scenario, maybe the user is acting as the negotiator and needs to generate a message that includes ARGUMENT.} \\
\\ 
\texttt{Wait, the user's message here is a simulation of the negotiation. The user is supposed to write a message that follows the protocol. The previous error was that the user's response didn't have an ARGUMENT. So the user needs to correct that.} \\
\\ 
\texttt{In the example given, the assistant's response included both PROPOSAL and ARGUMENT. So the user's next message should include both. Let me think of a good proposal and an argument.} \\
\\ 
\texttt{First, let's calculate the optimal set. Let's list all items with their effort and importance:} \\
\\ 
\texttt{Effort:} \\
\texttt{A19:767} \\
\texttt{B91:533} \\
\texttt{A30:514} \\
\texttt{B11:9} \\
\texttt{C47:539} \\
\texttt{C75:125} \\
\texttt{C26:153} \\
\texttt{C73:325} \\
\texttt{A47:937} \\
\texttt{A15:745} \\
\texttt{A75:334} \\
\texttt{A39:804} \\
\texttt{A56:336} \\
\texttt{C24:587} \\
\texttt{B26:71} \\
\\ 
\texttt{Importance:} \\
\texttt{A19:138} \\
\texttt{B91:583} \\
\texttt{A30:868} \\
\texttt{B11:822} \\
\texttt{C47:783} \\
\texttt{C75:65} \\
\texttt{C26:262} \\
\texttt{C73:121} \\
\texttt{A47:508} \\
\texttt{A15:780} \\
\texttt{A75:461} \\
\texttt{A39:484} \\
\texttt{A56:668} \\
\texttt{C24:389} \\
\texttt{B26:808} \\
\\ 
\texttt{First, note that B11 has the highest importance (822) and very low effort (9). Including B11 is a must. Similarly, B26 has high importance (808) and effort 71. These two are probably essential.} \\
\\ 
\texttt{Next, let's look for other high{-}importance items. B91 has importance 583, effort 533} \\
            }
        }
    }
    & & \\ \\

    \theutterance \stepcounter{utterance}  
    & & & \multicolumn{2}{p{0.3\linewidth}}{
        \cellcolor[rgb]{0.9,0.9,0.9}{
            \makecell[{{p{\linewidth}}}]{
                \texttt{\tiny{[GM$|$GM]}}
                \texttt{abort} \\
            }
        }
    }
    & & \\ \\

    \theutterance \stepcounter{utterance}  
    & & & \multicolumn{2}{p{0.3\linewidth}}{
        \cellcolor[rgb]{0.9,0.9,0.9}{
            \makecell[{{p{\linewidth}}}]{
                \texttt{\tiny{[GM$|$GM]}}
                \texttt{end game} \\
            }
        }
    }
    & & \\ \\

\end{supertabular}
}

\end{document}
